\documentclass[10pt]{article}
\usepackage[margin=1in]{geometry}
\usepackage{graphicx}
\usepackage{booktabs}
\usepackage{float}
\usepackage{longtable}
\usepackage{hyperref}
\usepackage{microtype}
\usepackage[numbers]{natbib}
\usepackage[font=small,labelfont=bf,skip=4pt]{caption}
\setlength{\parskip}{0pt}
\widowpenalty=10000
\clubpenalty=10000
\title{ Installed Ground Truth for Training-Data Attribution: Surface Form Shapes Causal Potency and Single-Checkpoint TracIn Removal Backfires }
\author{ Anonymous }
\date{}

% Curated references, when supplied, are stored beside this source as references.bib.




\begin{document}
\maketitle




\begin{abstract}
Training-data attribution is typically validated against content-similarity proxies rather than causal effects known in advance. We therefore install per-source effects by fine-tuning, measure each against a retrained control, remove each method’s top-ranked documents, and retrain. This exploratory study examines factory farming with Qwen3-4B, gradient-based removal methods, and paired training seeds for the declared seed-paired families; other readings are labeled separately. The same premise specification produced a 22x range of installed effects across surface forms, from +0.0072 for long articles at the two-epoch reading to +0.1574 for short answers, although the pair also differs in premise density and voice. Factorial cells were separately estimated corpus arms, and their ratios were point-estimate comparisons rather than tested effects. In the evaluated cells, oracle removal, a recorded positive-control contrast rather than a candidate method, recovered attributable fractions of +0.52 and +0.62 at the twenty-percent budget. By contrast, removal guided by the single-final-checkpoint TracIn implementation increased the effect at both seeds at the ten-percent budget, with attributable fractions of -0.91 and -0.42 and intervals excluding zero. These removability and detection inferences are cell-local.


\end{abstract}




\section{ Introduction }
\label{sec:introduction}

Training-data attribution is typically validated against content-similarity proxies rather than causal effects known in advance. Here, the ground truth is instead an installed, separately measured causal effect size.



This distinction matters because content alone did not determine causal potency: the same premise specification produced markedly different installed effects for a short-answer corpus and a long-article corpus at the two-epoch reading. Length, premise density, and voice all differ across this pair, and the long endpoint rises at the three-epoch reading. The separately estimated factorial corpus arms distinguish length from premise density: the short-dense cell was more potent than either the long-dense cell at matched density or the short-sparse cell at matched length. These are point-estimate ratios across separately estimated corpus arms, not tested effects.



The recorded oracle positive control further showed, in the evaluated cells, that the installed effect was removable and that the instrument detected removal; the oracle is a positive control, not a candidate, and the inference is restricted to those cells. By contrast, at the narrower budget, removing documents ranked highest by TracIn increased the installed effect at both seeds, with the two intervals excluding zero in the same counterproductive direction. This is a sign test against zero and does not depend on the overlap criterion. The result is scoped to the single-final-checkpoint implementation evaluated here.



The study is exploratory, covering one topic, factory farming; one model, Qwen3-4B; gradient-based removal methods; and paired training seeds for the declared seed-paired families. Two-seed replication applies only to declared seed-paired families. Other readings are labeled separately.





\section{Related Work}
\label{sec:related_work}
\paragraph{Training-data attribution and its evaluation.}
Influence-function and gradient-similarity methods estimate which training examples drive a
model behaviour \citep{koh2017influence,pruthi2020tracin,grosse2023influence}. Their evaluation
has largely relied on proxy ground truth. FTRACE scores attribution against entailment-labelled
fact tracing and already reports BM25 beating TracIn \citep{akyurek2022ftrace}; DATE-LM
benchmarks attribution across applications, including retrain-based validation of top-ranked
removal, and finds simple baselines competitive \citep{jiao2025datelm}; constructed-provenance
testbeds inject fabricated articles while explicitly disclaiming causal claims
\citep{li2025fakewiki}. The length confound our baseline embodies is likewise established:
RelatIF diagnoses norm-dominated rankings \citep{barshan2020relatif}, LESS calls the issue well
known and prescribes the cosine normalisation we also test \citep{xia2024less}, TrackStar meets
the same confound from the long-document end at pretraining scale \citep{chang2025trackstar},
and influence-function-family methods are reported fragile on large models
\citep{li2025influencework}. What this paper adds to that line is the \emph{kind} of ground
truth: prior testbeds label sources with binary provenance, acquisition, or entailment, whereas
here every source carries a separately measured causal effect size netted against a retrained
control -- so a method is scored against how much a source moved the model, and a source that
provably moved it by nothing is a planted trap rather than a positive. Two published positive
results bound our negative one: full-trajectory counterfactual influence separates provenance
where single-checkpoint scores struggle \citep{brazilek2026smalledits}, and TracIn's original
formulation sums over checkpoints along the trajectory \citep{pruthi2020tracin}; what fails
here is the single-checkpoint variant we test.

\paragraph{Installing beliefs by fine-tuning.}
Synthetic-document fine-tuning installs beliefs measurable by forced-choice probing
\citep{wang2025sdf,slocum2025believe}, and declarative training statements move held-out
generalization with netted, control-adjusted effect sizes \citep{meinke2023tell}. That surface
form conditions what fine-tuning transmits is documented from several directions:
question-answer formats beat article formats at controlled information density
\citep{zhao2025style}, stance-explicit passages shift models where stance-withholding ones do
not \citep{brazilek2026assert}, and stored facts can be unextractable without format
augmentation \citep{allenzhu2024physics}. Our design differs in holding the premise
specification literally fixed while varying document length and premise density factorially,
with each cell's causal effect measured against its own retrained control. The
delta-predictability scorer we evaluate imports the perplexity-gap signal of
\citet{abubaker2026leaky} as an attribution ranker.




\section{ Methods }
\label{sec:methods}

This exploratory study examined factory farming using Qwen3-4B, gradient-based removal methods, and paired training seeds for the declared seed-paired families. Two-seed replication applies only to declared seed-paired families; other readings are labeled separately. The pool comprised six source corpora about the same topic, sharing one premise specification while varying surface form by design, including an off-topic control and a source constructed to carry no causal effect.



Removal arms used full fine-tuning of Qwen3-4B under one recipe identical across arms; no removal-arm epoch count or learning rate is asserted, and no mechanism rationale is claimed. Dose was matched on training pairs and epochs: every arm retained the same row count and optimizer steps, so removed word mass varied by method. Each corpus was fine-tuned and evaluated with a forced-choice belief probe. The installed effect was the change in endorsement probability netted against a retrained matched off-topic control, because belief contrasts use matched off-topic-control netting. Before removal, this installed effect was +0.0167 at seed 42 and +0.0247 at seed 7, based on full fine-tuning and seed-specific readings.



Each method scored every pool document. The oracle, treated as a positive control, ranked sources by separately measured effect. In this implementation, TracIn scored gradient agreement at a single final checkpoint only; TracIn-cosine applied the prescribed length normalisation; delta-predictability ranked documents by the change in loss between the base and tuned model; and the word-count baseline removed the shortest documents first. The protocol removed the top-ranked pairs at budgets of ten and twenty percent of pool pairs, retrained from the base model under the identical recipe, and defined attributable fraction as AF = one minus the installed effect after removal divided by the installed effect before removal.



Intervals were estimated by a paired bootstrap over probe items with ten thousand draws, using ninety-five percent intervals throughout. Under the overlap criterion, a method failed to separate when its ninety-five percent interval overlapped the word-count baseline's interval in the same budget-and-seed cell. This criterion denotes interval overlap rather than a pairwise test; this caveat is stated only in methods.







\section{ Results }
\label{sec:results}

Across the four separately estimated corpus arms, factorial point estimates were +0.0072 for long-sparse, +0.0547 for long-dense, +0.0506 for short-sparse, and +0.1190 for short-dense. At the strong cells, the short-dense estimate was 2.2x the long-dense estimate at matched density and 2.4x the short-sparse estimate at matched length; these ratios are point-estimate comparisons rather than tested effects. Measured from the weakest cell, the corresponding ratios were 7.0x for length at sparse density and 7.6x for density at long length. This descriptive pattern was consistent with interaction between the factors, although no interaction test was computed.



The dose-matched descriptive-conclusions corpus had an installed effect of +0.1106; this was a dose-matched comparison point only.



Across six sources, the embedding retriever’s Spearman rank correlation with the installed ordering was +0.8857 at positive polarity and +0.9429 at negative polarity, compared with +0.2571 at both polarities for the length baseline. The retriever used the recorded chunked maximum-over-chunks score, and no removal was run for this retriever; ranking fidelity therefore was not tested as removal efficacy. The intervals were bootstrap-over-documents rank-stability bands and were quantized because rank correlation over six items is discrete.



At the twenty-percent budget, the oracle recovered +0.52 at seed 42 and +0.62 at seed 7, with intervals excluding zero at both seeds. As a positive-control contrast rather than a candidate, these cells showed that the installed effect was removable and that the instrument detected removal; inference was restricted to the declared twenty-percent cells.



At the ten-percent budget, the single-final-checkpoint TracIn implementation produced attributable fractions of -0.91 at seed 42 and -0.42 at seed 7, with both intervals excluding zero. Removing its top-ranked documents increased the installed effect; this sign test against zero was independent of the overlap criterion. TracIn-cosine yielded -0.95 at seed 42 with an interval excluding zero and -0.20 at seed 7 with an interval that straddled zero, so the interval-excluding counterproductive reading was single-seed for this method.



At the ten-percent budget, delta-predictability flipped sign across seeds: -0.58 at seed 42 and +0.34 at seed 7, with both intervals excluding zero. This seed contradiction, rather than an average, left the result unresolved.





\begin{table}[H]
\centering
\footnotesize
\caption{ Installed effect across the length by premise-density factorial corpus arms, with rows for document length and columns for premise density. Each cell reports the separately estimated effect with its 95\% interval. Source: paper/synthesis.json. }
\label{tab:factorial_2x2}


\resizebox{0.98\linewidth}{!}{

\begin{tabular}{ lll }
\toprule
document length & sparse premises (\textasciitilde{}4\%) & dense premises (\textasciitilde{}15\%) \\
\midrule

\textasciitilde{}101 words & \textbf{ +0.051 [+0.031, +0.072] } & \textbf{ +0.12 [+0.090, +0.15] } \\

\textasciitilde{}650--726 words & \textbf{ +0.0072 [+0.0016, +0.014] } & \textbf{ +0.055 [+0.039, +0.071] } \\

\bottomrule
\end{tabular}
}
\par\footnotesize\emph{Note.} Cells are separately estimated corpus arms with 95\% bootstrap intervals; no paired contrast between cells was computed. Bold excludes zero.
\end{table}




\begin{table}[tbp]
\centering
\footnotesize
\caption{ Attributable fraction and same-cell interval overlap with the word-count baseline by method, removal budget, and seed. The oracle overlaps the baseline at the ten-percent budget in both seeds, while the table carries the full overlap pattern. Source: paper/synthesis.json. }
\label{tab:af_overlap}



\begin{tabular}{ p{7.20cm}llll }
\toprule
method & budget & seed & attributable fraction & non-separation criterion \\
\midrule

\textbf{ Positive control: removal ranked by the separately measured effect } &  &  &  &  \\

oracle (measured effect) & 10\% & 42 & -0.10 [-0.95, +0.35] & overlaps \\

oracle (measured effect) & 10\% & 7 & +0.30 [-0.02, +0.55] & overlaps \\

oracle (measured effect) & 20\% & 42 & \textbf{ +0.52 [+0.17, +0.86] } & does NOT overlap \\

oracle (measured effect) & 20\% & 7 & \textbf{ +0.62 [+0.29, +0.84] } & overlaps \\

\textbf{ Candidate methods } &  &  &  &  \\

TracIn & 10\% & 42 & \textbf{ -0.91 [-2.15, -0.28] } & overlaps \\

TracIn & 10\% & 7 & \textbf{ -0.42 [-1.09, -0.02] } & overlaps \\

TracIn & 20\% & 42 & -0.47 [-1.51, +0.19] & overlaps \\

TracIn & 20\% & 7 & +0.17 [-0.38, +0.52] & overlaps \\

TracIn-cosine & 10\% & 42 & \textbf{ -0.95 [-2.32, -0.32] } & overlaps \\

TracIn-cosine & 10\% & 7 & -0.20 [-0.77, +0.16] & overlaps \\

TracIn-cosine & 20\% & 42 & -0.05 [-0.76, +0.33] & overlaps \\

TracIn-cosine & 20\% & 7 & +0.22 [-0.17, +0.48] & overlaps \\

delta-predictability & 10\% & 42 & \textbf{ -0.58 [-1.67, -0.02] } & overlaps \\

delta-predictability & 10\% & 7 & \textbf{ +0.34 [+0.08, +0.55] } & overlaps \\

delta-predictability & 20\% & 42 & +0.13 [-0.80, +0.59] & overlaps \\

delta-predictability & 20\% & 7 & +0.24 [-0.30, +0.62] & overlaps \\

\textbf{ Comparator: the word-count baseline itself } &  &  &  &  \\

word count (baseline) & 10\% & 42 & -0.49 [-1.53, +0.18] & (comparator) \\

word count (baseline) & 10\% & 7 & -0.14 [-0.68, +0.22] & (comparator) \\

word count (baseline) & 20\% & 42 & -0.30 [-0.94, +0.09] & (comparator) \\

word count (baseline) & 20\% & 7 & +0.12 [-0.37, +0.45] & (comparator) \\

\bottomrule
\end{tabular}

\par\footnotesize\emph{Note.} Verdict applies the pre-registered non-separation criterion: overlap of 95\% bootstrap intervals within a cell. Overlap is not a pairwise significance test; no pairwise contrast was computed. Bold excludes zero.
\end{table}


\begin{figure}[H]
\centering
\includegraphics[width=\linewidth]{ figures/attributable_fraction.png }
\caption{ Attributable fraction by method and removal budget, full fine-tuning, 95\% intervals. TracIn and TracIn-cosine sit left of zero at the ten-percent budget; only the twenty-percent oracle separates from zero.}
\label{fig:a_af}
\end{figure}






\section{ Discussion }
\label{sec:discussion}

At the narrow budget, the oracle and word-count intervals overlapped in both seed-specific cells despite their differing point-estimate directions. This overlap alone does not establish equivalence or any pairwise ordering. We interpret it as limited resolving power of the overlap criterion at these interval widths, so the narrow-budget verdict says little about method quality; this reading is specific to these cells and does not characterize the design’s overall diagnostic power.







\section{ Limitations }
\label{sec:limitations}

Replication is limited to two training seeds, so the source of the delta-predictability disagreement remains unresolved rather than being attributable to noise. The small denominator in the attributable fraction is a plausible but untested contributor to the wide intervals and sign changes across seeds, not an explanation of them.



In the long-article cell, the larger installed effect at the three-epoch reading than at the two-epoch reading indicates that the estimated gap is reading-step dependent; the direction also replicates at seed 7. This inference is local to that cell.



A stronger design would include a dose-matched random-removal arm retrained under the same full fine-tuning procedure to exclude the alternative that perturbing the corpus at all moves a small quantity. A separate stronger comparison would evaluate the present formulation against one that sums gradient agreement along the training trajectory; that comparison lies outside this evaluation.







\section{ Conclusion }
\label{sec:conclusion}

Across the factorial corpus arms evaluated, installed-effect point estimates ranged from +0.0072 in the long-sparse cell to +0.1190 in the short-dense cell. Because these corpus arms were estimated separately, this is a descriptive pattern rather than a tested main effect of either factor.



The positive-control contrast at the twenty-percent budget showed that removal ranked by separately measured effect eliminated roughly half or more of the installed effect at both seeds. This supports removability in the declared cells and confirms that the instrument detected removal there, without extending that interpretation beyond those cells.



By contrast, at the ten-percent budget, removal guided by the single-final-checkpoint TracIn implementation evaluated here increased the installed effect at both seeds, yielding negative attributable fractions in each case. This counterproductive direction is specific to the evaluated implementation and cells.







\appendix
\section{Supplementary diagnostics}




{\footnotesize
\begin{longtable}{ p{3.52cm}lp{8.16cm} }

\caption{ Full contrast ledger for all declared installed-effect, retrieval, pool, control, and attributable-fraction contrasts. Estimates, 95\% intervals, seed or reading scope, and recorded qualifications are shown for provenance and verification. Source: paper/synthesis.json. \emph{Note.} Rows without intervals are deterministic point contrasts over recorded arm scores. }
\label{tab:contribution_ladder}\\
\toprule
intervention & net effect & qualification \\
\midrule
\endfirsthead
\toprule
intervention & net effect & qualification \\
\midrule
\endhead
\bottomrule
\endfoot
\bottomrule
\endlastfoot

\textbf{ Installed effects } &  &  \\

Installed effect: descriptive conclusions (dose-matched comparison) & \textbf{ +0.11 [+0.087, +0.14] } & Conclusion-stating corpus at matched dose. \\

Installed effect: premises as \textasciitilde{}726-word articles (long-sparse cell) & \textbf{ +0.0072 [+0.0016, +0.014] } & The long-sparse cell of the 2x2 (\textasciitilde{}726 words, \textasciitilde{}4.0\% premise density, third-person articles), read at two epochs; matched off-topic control retrained per seed. See the step-36 row for the reading-step dependence. \\

Installed effect: the SAME premises as \textasciitilde{}104-word first-person answers & \textbf{ +0.16 [+0.13, +0.19] } & Identical premise specification to the long-article row; the point-estimate ratio between the two rows is 22x. Length, premise density, and voice all differ across this pair; the 2x2 cells separate the first two. \\

Installed effect: long documents, dense premises (long-dense cell) & \textbf{ +0.055 [+0.039, +0.071] } & \textasciitilde{}650 words at \textasciitilde{}14.9\% density. Measured from the weakest cell the single-factor ratios are larger (7.6x for density at long length, 7.0x for length at sparse density), so the two factors interact rather than compose multiplicatively. \\

Installed effect: the same long-article corpus read at three epochs & \textbf{ +0.034 [+0.025, +0.044] } & The long-sparse cell rises 4.75x between the two-epoch and three-epoch readings with non-overlapping intervals, so the size of the surface-form gap depends on the reading step; the rise replicates in direction at seed 7. \\

Installed effect: short answers, dense premises (short-dense cell) & \textbf{ +0.12 [+0.090, +0.15] } & Third person, \textasciitilde{}101 words, \textasciitilde{}14.9\% premise density. Its point estimate is 2.2x the long-dense cell's at matched density and 2.4x the short-sparse cell's at matched length. \\

Installed effect: short answers, sparse premises (short-sparse cell) & \textbf{ +0.051 [+0.031, +0.072] } & \textasciitilde{}101 words at \textasciitilde{}4.6\% density: same length as the short-dense cell, same density class as the long-article cell. \\

\textbf{ Pool and controls } &  &  \\

Pool effect before removal, seed 42 (dB\_before) & \textbf{ +0.017 [+0.0084, +0.026] } & Denominator of every AF at this seed; full fine-tuning. \\

Pool effect before removal, seed 7 (dB\_before) & \textbf{ +0.025 [+0.014, +0.037] } & Replicates seed 42 with overlapping intervals. \\

Off-topic control machinery, seed 42 & -0.00042 [-0.0034, +0.0029] & Straddles zero, as a clean control should. \\

\textbf{ Attributable fraction, by method } &  &  \\

AF, oracle removal at 10\% budget, seed 42 & -0.10 [-0.95, +0.35] & Straddles zero. Belongs to the power diagnostic; wide intervals here do not establish an ordering against the 20\% cells. \\

AF, oracle removal at 10\% budget, seed 7 & +0.30 [-0.016, +0.55] & Straddles zero; the power diagnostic's second cell. \\

AF, oracle removal at 20\% budget, seed 42 & \textbf{ +0.52 [+0.17, +0.86] } & Positive control: removal ranked by the separately measured per-source effect. \\

AF, oracle removal at 20\% budget, seed 7 & \textbf{ +0.62 [+0.29, +0.84] } & Seed-7 replication of the positive control; spread 0.10. \\

AF, delta-predictability at 10\% budget, seed 42 & \textbf{ -0.58 [-1.67, -0.018] } & Excludes zero NEGATIVE; contradicts the same cell at seed 7. Unresolved. \\

AF, delta-predictability at 10\% budget, seed 7 & \textbf{ +0.34 [+0.078, +0.55] } & Excludes zero POSITIVE; contradicts seed 42. The pair is a seed contradiction, not an average. \\

AF, delta-predictability at 20\% budget, seed 42 & +0.13 [-0.80, +0.59] & Straddles zero at the wider budget. \\

AF, delta-predictability at 20\% budget, seed 7 & +0.24 [-0.30, +0.62] & Straddles zero; the wider budget does not resolve the seed contradiction at 10\%. \\

AF, TracIn removal at 10\% budget, seed 42 & \textbf{ -0.91 [-2.15, -0.28] } & Negative and excluding zero: removing top-ranked documents INCREASED the effect. A sign test against zero, independent of the overlap criterion. \\

AF, TracIn removal at 10\% budget, seed 7 & \textbf{ -0.42 [-1.09, -0.022] } & Seed-7 replication of the counterproductive direction, interval excluding zero. \\

AF, TracIn removal at 20\% budget, seed 42 & -0.47 [-1.51, +0.19] & Straddles zero at the wider budget. \\

AF, TracIn removal at 20\% budget, seed 7 & +0.17 [-0.38, +0.52] & Straddles zero; completes the TracIn cell block. \\

AF, TracIn-cosine removal at 10\% budget, seed 42 & \textbf{ -0.95 [-2.32, -0.32] } & The prescribed length-normalisation fix; also counterproductive at this seed and budget, interval excluding zero. Single-seed for this method. \\

AF, TracIn-cosine at 10\% budget, seed 7 & -0.20 [-0.77, +0.16] & Straddles zero at this seed; the counterproductive reading is single-seed for this method. \\

AF, TracIn-cosine removal at 20\% budget, seed 42 & -0.053 [-0.76, +0.33] & Straddles zero at the wider budget. \\

AF, TracIn-cosine at 20\% budget, seed 7 & +0.22 [-0.17, +0.48] & Straddles zero; completes the TracIn-cosine cell block. \\

AF, word-count baseline at 10\% budget, seed 42 & -0.49 [-1.53, +0.18] & The model-free comparator. Non-separation is asserted as interval OVERLAP with each gradient-based method in the same budget-and-seed cell; no pairwise contrast was computed and none may be implied. \\

AF, word-count baseline at 10\% budget, seed 7 & -0.14 [-0.68, +0.22] & Comparator replication; the same interval-overlap criterion applies. \\

AF, word-count baseline at 20\% budget, seed 42 & -0.30 [-0.94, +0.087] & Comparator at the wider budget; the same interval-overlap criterion applies. \\

AF, word-count baseline at 20\% budget, seed 7 & +0.12 [-0.37, +0.45] & Completes the comparator block. \\

\textbf{ Retrieval rank correlations (not attributable fractions) } &  &  \\

Rank correlation with installed effect: embedding retriever, negative polarity & \textbf{ +0.94 [+0.89, +0.94] } & Same reading at the other polarity. A rank correlation over 6 sources; ranking fidelity is not removal efficacy. \\

Rank correlation with installed effect: embedding retriever, positive polarity & \textbf{ +0.89 [+0.89, +0.94] } & Spearman over 6 sources, not an attributable fraction; no removal was run for this retriever. intfloat/e5-base-v2, chunked, score = max over chunks; the mean-over-chunks variant gives +0.9429 at both polarities. Not length-confounded: rho(score, words) is -0.18 to -0.29. The interval is a bootstrap over documents within source -- a rank-stability band, quantized because a rank correlation over six items is discrete, so it can touch the point estimate. \\

Rank correlation with installed effect: length baseline, negative polarity & \textbf{ +0.26 [+0.14, +0.66] } & Baseline at the other polarity. Rank correlation only. \\

Rank correlation with installed effect: length baseline, positive polarity & \textbf{ +0.26 [+0.14, +0.43] } & The model-free baseline the retriever is read against, from the same bootstrap. Rank correlation only. \\

\end{longtable}
}




\begin{table}[tbp]
\centering
\footnotesize
\caption{ Run, artifact, quantity, and checkpoint behind every declared contrast. Source: paper/config.resolved.yaml. }
\label{tab:provenance}


\resizebox{0.98\linewidth}{!}{

\begin{tabular}{ lllll }
\toprule
contrast & run & artifact & quantity & checkpoint \\
\midrule

ladder\_evidence\_long & matrix\_v1\_step24 & belief\_summary.yaml & delta\_net & checkpoint-24 \\

ladder\_long\_step36 & matrix\_v1\_step36 & belief\_summary.yaml & delta\_net & checkpoint-36 \\

ladder\_evidence\_short & ms\_arms & belief\_summary.yaml & delta\_net & checkpoint-24 \\

ladder\_conclusions & matrix\_md\_2ep & belief\_summary.yaml & delta\_net & checkpoint-24 \\

ladder\_short\_dense & ms3p\_arms & belief\_summary.yaml & delta\_net & final/summary \\

ladder\_long\_dense & mld\_arms & belief\_summary.yaml & delta\_net & final/summary \\

ladder\_short\_sparse & ms\_sparse\_arms & belief\_summary.yaml & delta\_net & final/summary \\

rank\_e5\_positive & attrib\_mix\_v4 & retrieval\_summary.yaml & rho\_e5\_positive & final/summary \\

rank\_e5\_negative & attrib\_mix\_v4 & retrieval\_summary.yaml & rho\_e5\_negative & final/summary \\

rank\_neg\_length\_positive & attrib\_mix\_v4 & retrieval\_summary.yaml & rho\_neg\_length\_positive & final/summary \\

rank\_neg\_length\_negative & attrib\_mix\_v4 & retrieval\_summary.yaml & rho\_neg\_length\_negative & final/summary \\

pool\_db\_before\_s42 & attrib\_mix\_v4 & af\_summary.yaml & db\_before\_s42 & final/summary \\

pool\_db\_before\_s7 & attrib\_mix\_v4 & af\_summary.yaml & db\_before\_s7 & final/summary \\

pool\_machinery\_s42 & attrib\_mix\_v4 & af\_summary.yaml & machinery\_s42 & final/summary \\

af\_oracle\_p20\_s42 & attrib\_mix\_v4 & af\_summary.yaml & af\_oracle\_p20\_s42 & final/summary \\

af\_oracle\_p20\_s7 & attrib\_mix\_v4 & af\_summary.yaml & af\_oracle\_p20\_s7 & final/summary \\

af\_oracle\_p10\_s42 & attrib\_mix\_v4 & af\_summary.yaml & af\_oracle\_p10\_s42 & final/summary \\

af\_oracle\_p10\_s7 & attrib\_mix\_v4 & af\_summary.yaml & af\_oracle\_p10\_s7 & final/summary \\

af\_tracin\_p10\_s42 & attrib\_mix\_v4 & af\_summary.yaml & af\_tracin\_p10\_s42 & final/summary \\

af\_tracin\_p10\_s7 & attrib\_mix\_v4 & af\_summary.yaml & af\_tracin\_p10\_s7 & final/summary \\

af\_tracin\_p20\_s42 & attrib\_mix\_v4 & af\_summary.yaml & af\_tracin\_p20\_s42 & final/summary \\

af\_tracin\_p20\_s7 & attrib\_mix\_v4 & af\_summary.yaml & af\_tracin\_p20\_s7 & final/summary \\

af\_tracin\_cos\_p10\_s42 & attrib\_mix\_v4 & af\_summary.yaml & af\_tracin\_cos\_p10\_s42 & final/summary \\

af\_tracin\_cos\_p10\_s7 & attrib\_mix\_v4 & af\_summary.yaml & af\_tracin\_cos\_p10\_s7 & final/summary \\

af\_tracin\_cos\_p20\_s42 & attrib\_mix\_v4 & af\_summary.yaml & af\_tracin\_cos\_p20\_s42 & final/summary \\

af\_tracin\_cos\_p20\_s7 & attrib\_mix\_v4 & af\_summary.yaml & af\_tracin\_cos\_p20\_s7 & final/summary \\

af\_delta\_pred\_p10\_s42 & attrib\_mix\_v4 & af\_summary.yaml & af\_delta\_pred\_p10\_s42 & final/summary \\

af\_delta\_pred\_p10\_s7 & attrib\_mix\_v4 & af\_summary.yaml & af\_delta\_pred\_p10\_s7 & final/summary \\

af\_delta\_pred\_p20\_s42 & attrib\_mix\_v4 & af\_summary.yaml & af\_delta\_pred\_p20\_s42 & final/summary \\

af\_delta\_pred\_p20\_s7 & attrib\_mix\_v4 & af\_summary.yaml & af\_delta\_pred\_p20\_s7 & final/summary \\

af\_wordcount\_p10\_s42 & attrib\_mix\_v4 & af\_summary.yaml & af\_wordcount\_p10\_s42 & final/summary \\

af\_wordcount\_p10\_s7 & attrib\_mix\_v4 & af\_summary.yaml & af\_wordcount\_p10\_s7 & final/summary \\

af\_wordcount\_p20\_s42 & attrib\_mix\_v4 & af\_summary.yaml & af\_wordcount\_p20\_s42 & final/summary \\

af\_wordcount\_p20\_s7 & attrib\_mix\_v4 & af\_summary.yaml & af\_wordcount\_p20\_s7 & final/summary \\

\bottomrule
\end{tabular}
}
\par\footnotesize\emph{Note.} Training seeds: 42 (default) and 7 (quantities suffixed s7). The removal pool was scored at checkpoint-23 of attrib\_mix\_v4; matched off-topic controls come from m0\_multiform at checkpoint-24.
\end{table}




\clearpage
\bibliographystyle{plainnat}
\bibliography{references}



\end{document}
